\documentclass[aspectratio=169,11pt]{beamer}

% ============================================================
% THEME & COLOURS (matched to GDP PowerPoint)
% ============================================================
\usetheme{default}
\usecolortheme{default}
\setbeamertemplate{navigation symbols}{}

\definecolor{gdpPurple}{HTML}{633E88}
\definecolor{gdpDarkBlue}{HTML}{0C406D}
\definecolor{gdpNavy}{HTML}{091932}
\definecolor{gdpOrange}{HTML}{DD7A0F}
\definecolor{gdpSlate}{HTML}{3B4950}
\definecolor{gdpLightGray}{HTML}{E7E6E6}
\definecolor{gdpBurgundy}{HTML}{820E21}

\setbeamercolor{background canvas}{bg=white}
\setbeamercolor{normal text}{fg=gdpSlate}
\setbeamercolor{frametitle}{fg=gdpPurple,bg=white}
\setbeamercolor{framesubtitle}{fg=gdpDarkBlue,bg=white}
\setbeamercolor{title}{fg=gdpPurple}
\setbeamercolor{subtitle}{fg=gdpDarkBlue}
\setbeamercolor{author}{fg=gdpSlate}
\setbeamercolor{institute}{fg=gdpSlate}
\setbeamercolor{date}{fg=gdpSlate}
\setbeamercolor{item}{fg=gdpDarkBlue}
\setbeamercolor{subitem}{fg=gdpPurple}
\setbeamercolor{block title}{fg=white,bg=gdpDarkBlue}
\setbeamercolor{block body}{fg=gdpSlate,bg=gdpLightGray!30}
\setbeamercolor{block title alerted}{fg=white,bg=gdpBurgundy}
\setbeamercolor{block body alerted}{fg=gdpSlate,bg=gdpLightGray!30}
\setbeamercolor{footline}{fg=gdpSlate}
\setbeamercolor{headline}{fg=white,bg=gdpPurple}

\usepackage{helvet}
\renewcommand{\familydefault}{\sfdefault}
\usepackage{amsmath,amssymb,graphicx,booktabs,array}
\usepackage{pifont}
\newcommand{\cmark}{\ding{51}}
\newcommand{\xmark}{\ding{55}}

\setlength{\tabcolsep}{3pt}
\setbeamertemplate{itemize/enumerate body begin}{\small}
\setbeamertemplate{itemize/enumerate subbody begin}{\footnotesize}
\addtolength{\leftmargini}{-0.3cm}
\addtolength{\leftmarginii}{-0.2cm}

\setbeamertemplate{frametitle}{%
  \vspace{0.1cm}
  \textbf{\large\insertframetitle}
  \vspace{0.05cm}
  \par\textcolor{gdpDarkBlue}{\rule{\textwidth}{1.2pt}}
  \vspace{-0.3cm}
}

\setbeamertemplate{footline}{%
  \leavevmode%
  \hbox{%
    \begin{beamercolorbox}[wd=.333\paperwidth,ht=2.5ex,dp=1ex,leftskip=0.3cm]{footline}%
      \insertshortauthor
    \end{beamercolorbox}%
    \begin{beamercolorbox}[wd=.333\paperwidth,ht=2.5ex,dp=1ex,center]{footline}%
      \insertshorttitle
    \end{beamercolorbox}%
    \begin{beamercolorbox}[wd=.333\paperwidth,ht=2.5ex,dp=1ex,rightskip=0.3cm]{footline}%
      \insertframenumber{}
    \end{beamercolorbox}%
  }%
}

\setbeamertemplate{title page}{
  \vspace{1.5cm}
  \begin{center}
    {\Huge\bfseries\textcolor{gdpPurple}{\inserttitle}}\\[0.4cm]
    {\Large\textcolor{gdpDarkBlue}{\insertsubtitle}}\\[1.2cm]
    {\large\textcolor{gdpSlate}{\insertauthor}}\\[0.3cm]
    {\normalsize\textcolor{gdpSlate}{\insertinstitute}}\\[0.3cm]
    {\normalsize\textcolor{gdpSlate}{\insertdate}}
  \end{center}
}

\title{Feasibility / Desirability Study}
\subtitle{Justifying the Pivot from Acoustic Wave Computing to Chemical / Ionic Computation}
\author{S. Venugopal}
\institute{Cranfield University\hspace{0.3cm}|\hspace{0.3cm}Supervisors: Prof. Weisi Guo, Prof. Takis Tsoutsanis}
\date{23 June 2026}

\begin{document}

\begin{frame}[plain]
  \titlepage
\end{frame}

% ============================================================
\begin{frame}{Purpose of this deck}
  This deck is a working document. It collects all the numbers, arguments, and caveats needed to decide whether to continue with acoustic porous wave computing or pivot to a chemical / ionic continuous-computing direction.

  \vspace{0.15cm}
  \textbf{Two questions:}
  \begin{enumerate}
    \item Is the current acoustic direction physically and strategically viable?
    \item If not, what is the best replacement within the remaining MSc timeline?
  \end{enumerate}

  \vspace{0.15cm}
  \textcolor{gdpSlate}{\small Approach: first-principles energy and length-scale estimates, plus scored feasibility/desirability metrics for each candidate.}
\end{frame}

% ============================================================
\begin{frame}{Why the acoustic direction is now in doubt}
  During the progress presentation and subsequent review, three issues emerged:
  \begin{enumerate}
    \item \textbf{Length scale:} at 2 kHz the acoustic wavelength in air is 17 cm, forcing a device that is tens of centimetres long.
    \item \textbf{Energy scale:} a quiet 1 Pa wave in a 1 cm$^2$ area for 1 ms costs $\sim$0.24 nJ, orders of magnitude above digital.
    \item \textbf{Model validity:} the Zwikker-Kosten (ZK) model is invalid at 2 kHz for typical foam pore sizes; the Johnson-Champoux-Allard (JCA) model is needed.
  \end{enumerate}

  \vspace{0.15cm}
  The supervisor's feedback: wave-physics physical neural networks are already well-covered; the project should return to concentrations, reactions, and ion-channel-like dynamics.
\end{frame}

% ============================================================
\begin{frame}{Metrics for comparison}
  \small
  \begin{columns}[T]
    \begin{column}{0.32\textwidth}
      \textbf{Physical performance}
      \begin{itemize}
        \item Energy per operation (J/op)
        \item Operation delay (s)
        \item Energy-delay product (J$\cdot$s)
        \item Device volume / density
        \item Wavelength or pattern length scale
      \end{itemize}
    \end{column}
    \begin{column}{0.32\textwidth}
      \textbf{Implementation feasibility}
      \begin{itemize}
        \item Forward model maturity
        \item Parameter availability
        \item Numerical tractability
        \item Experimental feasibility
        \item Readout / I/O feasibility
      \end{itemize}
    \end{column}
    \begin{column}{0.32\textwidth}
      \textbf{Research desirability}
      \begin{itemize}
        \item Novelty / differentiation
        \item Supervisor alignment
        \item Thesis coherence
        \item Timeline risk
        \item Defensibility at viva
      \end{itemize}
    \end{column}
  \end{columns}
\end{frame}

% ============================================================
\begin{frame}{Reference baselines}
  \small
  \begin{center}
    \begin{tabular}{lccc}
      \toprule
      System & Energy per op & Delay & Notes \\
      \midrule
      Thermodynamic minimum & $3 \times 10^{-21}$ J & -- & Landauer limit \\
      Biological synapse & $10^{-16}$ to $10^{-15}$ J & $\sim$1 ms & Nature's reference \\
      Digital MAC (7 nm AI accelerator) & $10^{-15}$ to $10^{-12}$ J & $\sim$1 ns & Current silicon baseline \\
      Digital SRAM access & $\sim 10^{-15}$ J & $\sim$1 ns & Memory reference \\
      Human brain (whole-system) & $\sim 10^{-14}$ J & $\sim$10 ms & Includes support metabolism \\
      \bottomrule
    \end{tabular}
  \end{center}

  \vspace{0.15cm}
  \textbf{Rule of thumb:} any physical computing substrate above $10^{-12}$ J per operation needs a very strong justification to compete with digital silicon.
\end{frame}

% ============================================================
\begin{frame}{Candidate 1: Acoustic Porous Wave Computing -- Energy}
  \small
  For a 1 Pa RMS acoustic wave in air:
  \[
    I = \frac{p_{\text{rms}}^2}{\rho_0 c_0} = \frac{1^2}{1.225 \times 343} \approx 2.4\ \text{mW/m}^2
  \]

  For a 1 cm$^2$ cross-section and a 1 ms operation:
  \[
    E_{\text{op}} \approx I \, A \, t = 2.4\ \text{mW/m}^2 \times 10^{-4}\ \text{m}^2 \times 10^{-3}\ \text{s} \approx 2.4 \times 10^{-10}\ \text{J} = 0.24\ \text{nJ}
  \]

  \vspace{0.1cm}
  \textbf{Context:}
  \begin{itemize}
    \item 0.24 nJ is about $10^2$ to $10^5$ times worse than a digital MAC.
    \item 1 Pa is already quiet; a louder source makes it worse.
    \item Most acoustic energy simply propagates through; only a small fraction is redirected to probes.
  \end{itemize}
\end{frame}

% ============================================================
\begin{frame}{Candidate 1: Acoustic Porous Wave Computing -- Length Scale}
  \small
  At 2 kHz:
  \[
    \lambda = \frac{c_0}{f} = \frac{343}{2000} = 0.1715\ \text{m} = 17.15\ \text{cm}
  \]

  Simulation geometry from \texttt{fdtd\_zk\_mxv\_v2.py}:
  \begin{center}
    \footnotesize
    \begin{tabular}{lcc}
      \toprule
      Feature & Physical size & In wavelengths \\
      \midrule
      Full domain & 0.50 m & 2.9 $\lambda$ \\
      Design region width & 0.20 m & 1.2 $\lambda$ \\
      Each 3x3 design block in x & 0.067 m & 0.39 $\lambda$ \\
      Source-to-probe distance & 0.34 m & 2.0 $\lambda$ \\
      \bottomrule
    \end{tabular}
  \end{center}

  \vspace{0.1cm}
  \textbf{Issue:} a diffractive computer usually needs many wavelengths to accumulate useful phase shifts. Here the whole device is only a few wavelengths long, and each design block is nearly half a wavelength wide. The geometry is too compact for the chosen frequency.
\end{frame}

% ============================================================
\begin{frame}{Candidate 1: Acoustic Porous Wave Computing -- Model Issues}
  \small
  \textbf{ZK vs JCA regime}
  \[
    \delta_v = \sqrt{\frac{2\eta}{\rho_0 \omega}} \approx 48.5\ \mu\text{m}
    \qquad
    Wo = \Lambda \sqrt{\frac{\rho_0 \omega}{\eta}} \approx 2.9\ \text{for}\ \Lambda = 100\ \mu\text{m}
  \]
  ZK assumes $Wo \ll 1$. At 2 kHz with 100 $\mu$m pores we are in the JCA regime.

  \vspace{0.15cm}
  \textbf{Sound-speed discontinuity in \texttt{fdtd\_zk\_2d\_v2.py}}
  \[
    K_{\text{eff}} = \frac{p_0}{\phi}\ \text{in porous region},
    \qquad
    K_{\text{eff}} = \rho_0 c_0^2\ \text{in air}
  \]
  This gives $c_{\text{eff}} \approx 288$ m/s in the porous region vs. 343 m/s in air. The discontinuity creates artificial reflections at interfaces.

  \vspace{0.15cm}
  \textbf{Thin absorbing sponge}
  \[
    \text{SPONGE\_WIDTH} = 10\ \text{cells} = 0.04\ \text{m} \approx 0.23\lambda
  \]
  A good sponge should be about one wavelength thick; the current one will reflect energy back.
\end{frame}

% ============================================================
\begin{frame}{Candidate 1: Acoustic Porous Wave Computing -- Summary}
  \small
  \textbf{Feasibility}
  \begin{itemize}
    \item Forward model: ZK solver exists but is invalid; JCA solver not yet implemented.
    \item Parameter availability: good (porous acoustic properties are tabulated).
    \item Numerical tractability: moderate (FDTD works, optimisation is slow).
  \end{itemize}

  \vspace{0.1cm}
  \textbf{Desirability}
  \begin{itemize}
    \item Energy: 0.24 nJ per op is 2 to 5 orders of magnitude worse than digital.
    \item Size: device must be 10 cm to 1 m because of long wavelength.
    \item Speed: ms latency, kHz bandwidth.
    \item Novelty: low; wave PNNs are well-covered.
    \item Supervisor alignment: poor.
  \end{itemize}

  \vspace{0.1cm}
  \begin{block}{Verdict}
    Acoustic porous wave computing is a poor fit for the project goals and timeline.
  \end{block}
\end{frame}

% ============================================================
\begin{frame}{Candidate 2: Chemical Reaction Networks (CRNs) -- Physics}
  \small
  A CRN is a set of reactions such as:
  \[
    A + B \xrightarrow{k_1} C,
    \qquad
    C \xrightarrow{k_2} D
  \]

  Under mass-action kinetics:
  \[
    \frac{d[C]}{dt} = k_1 [A][B] - k_2 [C]
  \]

  \vspace{0.1cm}
  \textbf{Energy per molecular transformation}
  \[
    kT\ \text{at 300 K} \approx 4 \times 10^{-21}\ \text{J}
    \qquad
    E_{\text{op}} \approx N \times kT
  \]

  \begin{center}
    \footnotesize
    \begin{tabular}{cc}
      \toprule
      Molecules per operation & Energy per operation \\
      \midrule
      1 (single-molecule limit) & $\sim$4 zJ \\
      $10^3$ & $\sim$4 aJ \\
      $10^6$ & $\sim$4 fJ \\
      $10^9$ & $\sim$4 pJ \\
      \bottomrule
    \end{tabular}
  \end{center}
\end{frame}

% ============================================================
\begin{frame}{Candidate 2: Chemical Reaction Networks (CRNs) -- Mapping to NNs}
  \small
  \textbf{How a CRN can compute:}
  \begin{itemize}
    \item Input: initial concentration of species $X_1, X_2, \dots$.
    \item Hidden state: concentrations of intermediate species during evolution.
    \item Output: steady-state concentrations of readout species $Y_1, Y_2, \dots$.
    \item Trainable parameters: rate constants $k_i$ and network topology.
  \end{itemize}

  \vspace{0.15cm}
  \textbf{Why this maps to a neural network}
  \begin{itemize}
    \item Mass-action ODEs are a continuous, nonlinear dynamical system.
    \item The stoichiometry and rate constants define a fixed computational graph.
    \item Supervised learning: choose $k_i$ so that target outputs match labels for training inputs.
    \item Recent reference: Nagipogu \& Reif (2025), ``Neural CRNs'' in ACS Synthetic Biology.
  \end{itemize}

  \vspace{0.15cm}
  \begin{block}{Key advantage}
    The forward model is a small system of ODEs; training uses standard SciPy + gradient descent.
  \end{block}
\end{frame}

% ============================================================
\begin{frame}{Candidate 2: Chemical Reaction Networks (CRNs) -- Summary}
  \small
  \textbf{Feasibility}
  \begin{itemize}
    \item Forward model: mature (mass-action ODEs).
    \item Parameter availability: excellent (rate constants are the design variables).
    \item Numerical tractability: high (scales to 10s--100s of species on a laptop).
    \item Experimental feasibility: moderate (DNA strand displacement, microfluidics, or well-mixed chemostats).
    \item Readout: fluorescence, spectroscopy, or electrochemical probes.
  \end{itemize}

  \vspace{0.1cm}
  \textbf{Desirability}
  \begin{itemize}
    \item Energy: can approach aJ per operation if $N$ is small.
    \item Size: microfluidic chambers or even single droplets.
    \item Speed: ms to s, depending on rates.
    \item Novelty: moderate; trainable CRNs are known but not saturated.
    \item Supervisor alignment: high (concentrations, reactions, ion-channel analogy).
  \end{itemize}
\end{frame}

% ============================================================
\begin{frame}{Candidate 3: Reaction-Diffusion Computing -- Physics}
  \small
  Governed by reaction-diffusion PDEs:
  \[
    \frac{\partial u}{\partial t} = D_u \nabla^2 u + f(u,v)
    \qquad
    \frac{\partial v}{\partial t} = D_v \nabla^2 v + g(u,v)
  \]

  \vspace{0.1cm}
  \textbf{Pattern-based computation}
  \begin{itemize}
    \item Inputs: initial concentrations or boundary stimuli.
    \item Computation: spatiotemporal pattern formation (waves, spots, Turing patterns).
    \item Outputs: presence/absence of a pattern at a probe location.
  \end{itemize}

  \vspace{0.1cm}
  \textbf{Energy and length scales}
  \[
    E_{\text{update}} \sim 2\ \text{nJ per mm}^3
    \qquad
    \lambda_{\text{pattern}} \approx \sqrt{D\tau} \sim 10\text{--}100\ \mu\text{m}
  \]
  \begin{itemize}
    \item Closer to spatial computing; good if supervisor insists on pattern dynamics.
    \item Harder to train because of PDE forward model and pattern sensitivity.
  \end{itemize}
\end{frame}

% ============================================================
\begin{frame}{Candidate 3: Reaction-Diffusion Computing -- Summary}
  \small
  \textbf{Feasibility}
  \begin{itemize}
    \item Forward model: PDE solver needed (finite-difference or FEniCS).
    \item Parameter availability: good for classical systems (Belousov-Zhabotinsky, Gray-Scott).
    \item Numerical tractability: moderate (2D time stepping is slower than ODEs).
    \item Experimental feasibility: moderate (BZ gel, microfluidic reactors).
    \item Readout: imaging or local electrode.
  \end{itemize}

  \vspace{0.1cm}
  \textbf{Desirability}
  \begin{itemize}
    \item Energy: competitive at the molecular scale.
    \item Size: mm-scale chips possible.
    \item Speed: seconds to minutes.
    \item Novelty: moderate; RD computing is well-known but trainable RD is less common.
    \item Supervisor alignment: moderate (matches ``pattern'' and ``spatial'' language).
  \end{itemize}
\end{frame}

% ============================================================
\begin{frame}{Candidate 4: Ion Channels / Excitable Membranes -- Physics}
  \small
  Hodgkin-Huxley style kinetics:
  \[
    C_m \frac{dV}{dt} = -\sum_i g_i (V - E_i)
  \]
  \[
    \frac{dg_i}{dt} = \alpha_i(V)(1 - g_i) - \beta_i(V)g_i
  \]

  \vspace{0.1cm}
  \textbf{Energy scales}
  \[
    E_{\text{ion}} = qV \approx 1.6 \times 10^{-20}\ \text{J}
    \qquad
    E_{\text{pump}} \sim 10 kT \approx 4 \times 10^{-20}\ \text{J per ion}
  \]
  For $10^6$ ions: $E_{\text{op}} \sim 40$ fJ.

  \vspace{0.1cm}
  \textbf{Length / time scales}
  \begin{itemize}
    \item Channel size: $\sim$10 nm; membrane thickness: $\sim$5 nm.
    \item Channel gating: $\mu$s; action-potential propagation: 0.1 to 100 m/s.
    \item Device: $\mu$m to mm.
  \end{itemize}
\end{frame}

% ============================================================
\begin{frame}{Candidate 4: Ion Channels / Excitable Membranes -- Summary}
  \small
  \textbf{Feasibility}
  \begin{itemize}
    \item Forward model: Hodgkin-Huxley or Markov-state ODEs are mature.
    \item Parameter availability: good for known channels; custom channels need characterisation.
    \item Numerical tractability: high (small ODE systems).
    \item Experimental feasibility: low for a custom thesis device (lipid bilayers, patch-clamp).
    \item Readout: electrical (pA currents) or fluorescence.
  \end{itemize}

  \vspace{0.1cm}
  \textbf{Desirability}
  \begin{itemize}
    \item Energy: extremely low per ion.
    \item Size: nanoscale channels.
    \item Speed: $\mu$s to ms.
    \item Novelty: high if framed as neuromorphic ionic computing.
    \item Supervisor alignment: high (explicitly requested ion-channel dynamics).
  \end{itemize}
\end{frame}

% ============================================================
\begin{frame}{Candidate 5: Organic / Iontronic Devices -- Physics}
  \small
  Examples: ion-gel transistors, electrochemical memristors, conducting-polymer devices.

  \vspace{0.15cm}
  \textbf{Governing physics}
  \[
    \sigma(x,t) = \sigma_{\text{max}}\, g(x,t)
  \]
  \[
    \frac{\partial g}{\partial t} = f(\phi, c_{\text{ion}}, g)
  \]

  \vspace{0.15cm}
  \textbf{Energy estimate}
  \[
    E = \tfrac{1}{2} C V^2 = 0.5 \times 1\ \text{fF} \times (0.1\ \text{V})^2 = 5 \times 10^{-18}\ \text{J} = 5\ \text{aJ (ideal)}
  \]
  Realistic devices: fJ to pJ per switching event.

  \vspace{0.1cm}
  \textbf{Why it is attractive}
  \begin{itemize}
    \item Direct bridge to electronics / neuromorphic hardware.
    \item Ion migration gives intrinsic memory (hysteresis).
  \end{itemize}
\end{frame}

% ============================================================
\begin{frame}{Candidate 5: Organic / Iontronic Devices -- Summary}
  \small
  \textbf{Feasibility}
  \begin{itemize}
    \item Forward model: coupled electrochemistry, ion transport, polymer physics.
    \item Parameters are device-specific and scarce.
    \item Drift-diffusion-Poisson plus hysteresis is complex to implement.
  \end{itemize}

  \vspace{0.1cm}
  \textbf{Desirability}
  \begin{itemize}
    \item Energy: very low per switching event.
    \item Size: thin-film devices, $\mu$m features.
    \item Speed: kHz to MHz, limited by ion mobility.
    \item Novelty: high; strong links to emerging neuromorphic iontronic literature.
    \item Supervisor alignment: high (ion dynamics).
  \end{itemize}

  \vspace{0.1cm}
  \begin{block}{Caveat}
    Best as a future direction; too device-specific and parameter-scarce for the remaining MSc timeline.
  \end{block}
\end{frame}

% ============================================================
\begin{frame}{Summary Scores}
  \small
  Score 1 (poor) to 5 (excellent). Mean rounded to one decimal.
  \begin{center}
    \footnotesize
    \begin{tabular}{lcccccccc}
      \toprule
      Candidate & Energy & Speed & Size & Feas. & Novelty & Fit & Risk & \textbf{Mean} \\
      \midrule
      Acoustic porous & 1 & 2 & 1 & 3 & 2 & 1 & 4 & \textbf{1.7} \\
      Reaction-diffusion & 2 & 1 & 3 & 3 & 3 & 4 & 3 & \textbf{2.7} \\
      Organic / iontronic & 4 & 3 & 4 & 2 & 5 & 3 & 1 & \textbf{3.2} \\
      Ion channels & 5 & 3 & 4 & 4 & 3 & 4 & 3 & \textbf{3.7} \\
      \textbf{CRNs} & 4 & 2 & 4 & 5 & 3 & 5 & 5 & \textbf{4.0} \\
      \bottomrule
    \end{tabular}
  \end{center}

  \vspace{0.15cm}
  \textbf{Interpretation:}
  \begin{itemize}
    \item CRNs win on feasibility, supervisor fit, and timeline risk.
    \item Ion channels are close but experimentally harder to realise within the MSc.
    \item Acoustic porous is the weakest option on nearly every metric.
  \end{itemize}
\end{frame}

% ============================================================
\begin{frame}{What is salvageable from the acoustic work?}
  \small
  The acoustic month is not wasted. The following pieces transfer directly to a CRN project:
  \begin{itemize}
    \item \textbf{Design-parameter abstraction:} mapping a discrete design grid to continuous physical parameters.
    \item \textbf{Forward-solver discipline:} time integration, CFL/stability checks, boundary-condition handling.
    \item \textbf{Optimisation loop:} grid search / random search / gradient-based tuning of design variables.
    \item \textbf{Readout abstraction:} probe concentrations (or pressures) at fixed locations and times.
    \item \textbf{Loss-function design:} classification / regression objective with regularisation.
  \end{itemize}

  \vspace{0.15cm}
  \textbf{What gets discarded:}
  \begin{itemize}
    \item The ZK / JCA porous acoustic model.
    \item The 2 kHz source frequency and cm-scale domain.
    \item The interpretation of the device as a diffractive acoustic computer.
  \end{itemize}
\end{frame}

% ============================================================
\begin{frame}{Recommendation}
  \small
  \begin{block}{Recommended direction}
    Pivot the thesis to \textbf{trainable Chemical Reaction Networks (CRNs)} as a continuous analog computing substrate.
  \end{block}

  \vspace{0.15cm}
  \textbf{Why CRNs?}
  \begin{itemize}
    \item Feasibility score 5/5: mass-action ODE forward model is immediate to implement.
    \item Supervisor fit 5/5: concentrations, reactions, and rate-gating map naturally to ion-channel-like dynamics.
    \item Timeline risk 5/5: first benchmark (XOR) can run within days; classifier benchmark within 1--2 weeks.
    \item Energy story is compelling: aJ to fJ per operation for modest molecule counts.
    \item Novelty is sufficient: trainable CRNs for neural computation are known but not saturated in the thesis literature.
  \end{itemize}

  \vspace{0.15cm}
  \textbf{Framing:} a ``single-pot'' continuous analog computer. Ion channels become the physical interpretation; CRNs become the tractable model system.
\end{frame}

% ============================================================
\begin{frame}{Proposed Timeline}
  \small
  \begin{center}
    \footnotesize
    \begin{tabular}{lll}
      \toprule
      Week & Task & Deliverable \\
      \midrule
      23--26 June & Lock direction with supervisor & Approved pivot \\
      29 June -- 3 July & CRN forward solver + XOR benchmark & Working Python code \\
      6--10 July & Trainable 2x2 MxV / small classifier & Optimised CRN \\
      13--17 July & Optional reaction-diffusion extension & Spatial logic gate \\
      20--24 July & Compare to acoustic baseline + write-up & Methodology chapter \\
      27--31 July & Final simulations + sensitivity study & Results chapter draft \\
      3--7 Aug & Write conclusions + final report & Submission-ready thesis \\
      \bottomrule
    \end{tabular}
  \end{center}

  \vspace{0.15cm}
  \textbf{Key assumption:} report chapters are modular; methodology and results can be rewritten without changing the introduction's framing of continuous analog computation.
\end{frame}

% ============================================================
\begin{frame}{Risk Mitigation}
  \small
  \textbf{Risk: CRNs are too abstract / not spatial enough.}
  \begin{itemize}
    \item Mitigation: frame CRNs as a ``single-pot'' continuous analog computer; add reaction-diffusion extension if needed.
  \end{itemize}

  \vspace{0.1cm}
  \textbf{Risk: Optimisation does not converge.}
  \begin{itemize}
    \item Mitigation: start with tiny benchmarks (XOR, 2x2 MxV) where convergence is easier; scale up gradually.
  \end{itemize}

  \vspace{0.1cm}
  \textbf{Risk: Supervisor wants ion channels specifically.}
  \begin{itemize}
    \item Mitigation: present CRNs as the tractable first step; ion channels become the aspirational physical interpretation (gated reaction rates).
  \end{itemize}

  \vspace{0.1cm}
  \textbf{Risk: Not enough time to rewrite report.}
  \begin{itemize}
    \item Mitigation: lock direction by end of June; keep report chapters modular so the methodology and results can be swapped without rewriting the entire thesis.
  \end{itemize}
\end{frame}

\end{document}
